\documentclass[11pt]{article}
\usepackage[T1]{fontenc}
\usepackage{lmodern}
\usepackage[margin=1in]{geometry}
\usepackage{amsmath,amssymb,amsthm,mathtools}
\usepackage{microtype}
\usepackage[colorlinks=true,linkcolor=blue,citecolor=blue,urlcolor=blue]{hyperref}
\newtheorem{theorem}{Theorem}[section]
\newtheorem{corollary}[theorem]{Corollary}
\newtheorem{lemma}[theorem]{Lemma}
\newtheorem{remark}[theorem]{Remark}
\DeclareMathOperator{\Res}{Res}
\newcommand{\Z}{\mathbb Z}
\newcommand{\Q}{\mathbb Q}
\title{A monic resultant for Bessel derivative coefficients}
\author{Keio}
\date{22 October 2026}
\begin{document}
\maketitle

\begin{abstract}
For the differential equation
$f''+cx^{-1}f'+(1+d(c-d-1)x^{-2})f=0$, we express each derivative
$f^{(n)}$ in the basis $(f,f')$ using two polynomial coefficients.
Their resultant with respect to $x$ is a monic polynomial in $c$ over
$\Z[d]$, of degree $n(n-1)$, for every $n\ge2$.
The specialization $d=0$, $c=2\nu+1$ proves the degree assertion
formulated in Section 2.3(a) of Lorch and Muldoon (1995) for derivatives
of $x^{-\nu}J_\nu(x)$. A classical theorem of Siegel then implies
that all nonzero zeros of these derivatives are transcendental whenever
$\nu\in\Q\setminus\tfrac12\Z$. More generally, the same conclusion
holds for $x^{-\mu}J_\nu(x)$ when $\mu-\nu$ or $\mu+\nu$ is an
algebraic integer. The remaining rational-root classification and the
all-orders question for derivatives of the unnormalized function
$J_\nu$ are not resolved here.
\end{abstract}

\section{The coefficient resultant}

Lorch and Muldoon \cite[Section 2.3]{LM95} studied the algebraic zeros
of successive derivatives of the normalized Bessel function
$x^{-\nu}J_\nu(x)$. Their polynomial elimination procedure led to
three conjectures. The first concerns only the degree of a resultant:
they proposed that its degree in $\nu$ is $n(n-1)$ for derivative
order $n$. We prove this assertion by identifying the highest-weight
coefficient polynomials, and obtain the stronger fact that a suitable
affine change of parameter makes the resultant monic.

The proof is algebraic. The only transcendence input, used subsequently
for zeros, is the classical nonvanishing and logarithmic-derivative
theorem for Bessel functions of rational order. In particular, no new
algebraic-independence theorem is required for the resultant calculation.

Let $c,d$ be indeterminates, and write
\[
 \Delta=d(c-d-1).
\]
Consider formally the second-order equation
\begin{equation}\label{eq:ode}
 f''+\frac c x f'+\left(1+\frac{\Delta}{x^2}\right)f=0.
\end{equation}
Define $A_n,B_n\in\Z[d,c,x]$ for $n\ge1$ by
\begin{equation}\label{eq:recurrence}
\begin{split}
 A_1&=0,\qquad B_1=1,\\
 A_{n+1}&=x\partial_xA_n-nA_n-(x^2+\Delta)B_n,\\
 B_{n+1}&=A_n+x\partial_xB_n+(1-c-n)B_n.
\end{split}
\end{equation}
Differentiating and substituting \eqref{eq:ode} gives
\begin{equation}\label{eq:derivative}
 f^{(n)}=x^{-n}\bigl(A_nf+xB_nf'\bigr)
\end{equation}
for every local solution on a domain avoiding zero. Both coefficient
polynomials are even in $x$.

\begin{theorem}\label{thm:main}
For every integer $n\ge2$, the polynomial
\[
 R_n(c,d)=\Res_x\bigl(A_n(x,c,d),B_n(x,c,d)\bigr)
\]
belongs to $\Z[d][c]$, is monic as a polynomial in $c$, and has
$c$-degree $n(n-1)$.
\end{theorem}

Our convention is the usual resultant with respect to $x$, without
removing factors corresponding to $x=0$. Removing such factors could
change the degree under discussion. At $d=0$, the polynomials denoted
by $p_n,q_n$ in \cite{LM95} satisfy
\[
 p_n=-B_n,\qquad q_n=-A_n.
\]
Their $x$-degrees are even, so sign changes and interchanging the two
arguments leave the resultant unchanged. Consequently Theorem
\ref{thm:main} proves their stated degree conjecture in exactly their
resultant convention. Since $c=2\nu+1$, the leading coefficient in
their original parameter $\nu$ is $2^{n(n-1)}$.

\section{The highest-weight calculation}

Treat $d$ as a coefficient, and assign weight one to both $x$ and $c$.
The additional potential $\Delta$ has weight at most one. The recurrences
therefore give
\[
 \deg A_n\le n,\qquad \deg B_n\le n-1,
\]
where these are total degrees in $x,c$. For $n\ge2$, the exact
$x$-degrees are
\begin{equation}\label{eq:x-degrees}
 d_A=2\lfloor n/2\rfloor,\qquad
 d_B=2\lfloor(n-1)/2\rfloor.
\end{equation}
We include the verification, as preservation of these degrees is needed
when reducing a resultant modulo a prime.

Let $a_n,b_n$ be the homogeneous parts of weights $n,n-1$ of $A_n,B_n$,
respectively. Terms involving $\Delta$ have lower weight, and hence
\begin{equation}\label{eq:top}
 a_n=-x^2b_{n-1},\qquad
 b_n=-cb_{n-1}-x^2b_{n-2}\quad(n\ge3),
 \qquad b_1=1,\quad b_2=-c.
\end{equation}
In particular, these highest-weight parts are independent of $d$.
The recurrence, or induction using Pascal's identity, gives
\begin{equation}\label{eq:binomial}
 b_n(x,c)=
 \sum_{j=0}^{\lfloor(n-1)/2\rfloor}
 (-1)^{n-1-j}\binom{n-1-j}{j}c^{n-1-2j}x^{2j}.
\end{equation}
The leading terms in this formula and $a_n=-x^2b_{n-1}$ give the
degrees in \eqref{eq:x-degrees}. The full recurrences imply the matching
upper bounds on $x$-degree, so there are no further terms of higher
$x$-degree in $A_n$ or $B_n$.

\begin{lemma}\label{lem:unit-resultant}
For every $n\ge2$,
\[
 \Res_x\bigl(a_n(x,1),b_n(x,1)\bigr)=1.
\]
\end{lemma}

\begin{proof}
For $n=2$, the pair is $(-x^2,-1)$ and the assertion is immediate.
Suppose $n\ge3$. We first check the absence of common finite zeros
over an arbitrary field, including fields of positive characteristic.
Since $b_n(0,1)=(-1)^{n-1}$, a common zero $\xi$ cannot be zero.
The equation $a_n(\xi,1)=0$ then gives $b_{n-1}(\xi,1)=0$.
The recurrence \eqref{eq:top}, using $\xi\ne0$, next gives
$b_{n-2}(\xi,1)=0$. Continuing backwards eventually contradicts
$b_1=1$.

It remains to handle the possible loss of degree after reduction modulo
a prime. If $n$ is even, the coefficient of the highest power of $x$
in $a_n(x,1)$ is $\pm1$. If $n$ is odd, the corresponding coefficient
in $b_n(x,1)$ is $\pm1$. These statements follow directly from
\eqref{eq:binomial}. Thus at least one of the two prescribed degrees
$d_A,d_B$ is preserved modulo every prime. Homogenize at precisely
these degrees. The homogenizations have no common zero at infinity,
and the preceding argument rules out a common finite zero. Their
resultant is consequently nonzero modulo every prime.

The integer $\Res_x(a_n(x,1),b_n(x,1))$ therefore has no prime divisor,
so it is $1$ or $-1$. For polynomials in $x^2$, the identity
\[
 \Res_x\bigl(P(x^2),Q(x^2)\bigr)
   =\Res_z(P(z),Q(z))^2
\]
shows that this integer is nonnegative. It must equal $1$.
\end{proof}

\begin{proof}[Proof of Theorem \ref{thm:main}]
Introduce a new indeterminate $T$, substitute $x=TX,c=T$, and divide
the two polynomials by their respective highest weights. In the
polynomial ring in $X$ with coefficients in $\Z[d][T^{-1}]$, we have
\[
 \begin{split}
 T^{-n}A_n(TX,T,d)&=a_n(X,1)+O(T^{-1}),\\
 T^{-(n-1)}B_n(TX,T,d)&=b_n(X,1)+O(T^{-1}).
 \end{split}
\]
Here $O(T^{-1})$ means a polynomial in $T^{-1}$ with zero constant
term. These expressions have prescribed $X$-degrees $d_A,d_B$.
The scaling laws for the resultant, followed by Lemma
\ref{lem:unit-resultant}, give
\[
 R_n(T,d)
 =T^{n d_B+(n-1)d_A-d_Ad_B}\bigl(1+O(T^{-1})\bigr).
\]
For $n=2m$ the exponent is $2m(2m-1)$, while for $n=2m+1$ it is
$2m(2m+1)$. Thus it equals $n(n-1)$ in both cases. Since the
resultant is in $\Z[d][T]$, its leading coefficient is exactly $1$.
\end{proof}

\section{Arithmetic consequences for zeros}

For rational $\nu$, the classical theorem of Siegel implies that
$J_\nu(\alpha)\ne0$ and $J_\nu'(\alpha)/J_\nu(\alpha)$ is
transcendental whenever $\alpha$ is nonzero and algebraic; see
\cite[Chapter 6, Section 2]{Shid89}, or the formulation used in
\cite{LM95}. For algebraic $\mu$, a local branch of
$f=x^{-\mu}J_\nu$ therefore satisfies
\begin{equation}\label{eq:siegel-input}
 f(\alpha)\ne0,\qquad
 \frac{f'(\alpha)}{f(\alpha)}
 =\frac{J_\nu'(\alpha)}{J_\nu(\alpha)}-\frac\mu\alpha
 \notin\overline{\Q}.
\end{equation}
This input is independent of the branch used near $\alpha$.

Put $d=\mu-\nu$ and $c=2\mu+1$. The transformed Bessel equation is
exactly \eqref{eq:ode}, since
\[
 \mu^2-\nu^2=d(c-d-1).
\]
If $f^{(n)}(\alpha)=0$ at a nonzero algebraic point, then
\eqref{eq:derivative} and \eqref{eq:siegel-input} imply
\begin{equation}\label{eq:common}
 A_n(\alpha,c,d)=B_n(\alpha,c,d)=0.
\end{equation}
Indeed, a nonzero coefficient $B_n$ would express $f'/f$ as an
algebraic number; after $B_n=0$, the nonvanishing of $f$ forces
$A_n=0$. In particular, \eqref{eq:common} gives $R_n(c,d)=0$.

\begin{corollary}\label{cor:normalized}
Let $\nu\in\Q\setminus\tfrac12\Z$. For every integer $n\ge0$,
every nonzero zero of
\[
 \left(\frac d{dx}\right)^n\bigl(x^{-\nu}J_\nu(x)\bigr)
\]
is transcendental.
\end{corollary}

\begin{proof}
The cases $n=0,1$ follow directly from \eqref{eq:siegel-input}.
For $n\ge2$, an algebraic exceptional point would make the rational
number $c=2\nu+1$ a root of the monic integer polynomial $R_n(c,0)$.
The rational-root theorem would then imply $c\in\Z$, contrary to
$2\nu\notin\Z$.
\end{proof}

\begin{corollary}[Integral shifts]\label{cor:shifts}
Let $\nu\in\Q\setminus\tfrac12\Z$ and
$\mu\in\overline{\Q}$. Suppose that at least one of $\mu-\nu$
and $\mu+\nu$ is an algebraic integer. For every $n\ge0$, every
nonzero zero of $(d/dx)^n(x^{-\mu}J_\nu(x))$ is transcendental.
\end{corollary}

\begin{proof}
Suppose first that $d=\mu-\nu$ is an algebraic integer. For $n\ge2$,
an exceptional point would give $R_n(c,d)=0$. Monicity makes $c$
integral over $\Z[d]$, hence integral over $\Z$. The number
$2\nu=c-2d-1$ would be a rational algebraic integer and therefore an
integer, a contradiction. If $\mu+\nu$ is integral instead, the same
argument uses the formal order $-\nu$ in the differential equation;
the chosen solution and the input \eqref{eq:siegel-input} are unchanged.
The cases $n=0,1$ again follow directly from that input.
\end{proof}

\begin{corollary}[Denominator obstruction]\label{cor:denominators}
Let $\mu,\nu\in\Q$. If a derivative of $x^{-\mu}J_\nu(x)$ has
a nonzero algebraic zero, every prime in the denominator of $2\nu$
occurs in the denominator of $\mu-\nu$ and also in the denominator
of $\mu+\nu$.
\end{corollary}

\begin{proof}
The cases of derivative order zero or one have no exceptional point
by \eqref{eq:siegel-input}, so suppose $n\ge2$.
If $d=\mu-\nu$ is integral at a prime $p$, then the coefficients
of the monic polynomial $R_n(c,d)$ lie in $\Z_{(p)}$. The rational
root $c$ is integral over this integrally closed ring, so $c$ and
$2\nu=c-2d-1$ are $p$-integral. Replacing the formal order by
$-\nu$ gives the other assertion.
\end{proof}

\section{Examples and the remaining questions}

The first two resultants provide short checks on the normalization.
With $\Delta=d(c-d-1)$, the recurrence gives
\[
 A_2=-x^2-\Delta,\qquad B_2=-c,
 \qquad R_2(c,d)=c^2.
\]
At the next order,
\[
 A_3=cx^2+(c+2)\Delta,\qquad
 B_3=c(c+1)-\Delta-x^2,
\]
and hence
\[
 R_3(c,d)=\bigl(c^2(c+1)+2\Delta\bigr)^2.
\]
Their degrees in $c$ are $2$ and $6$, respectively, as required.
For $d=0$, these become $c^2$ and $c^4(c+1)^2$.

The passage from common roots to a resultant has an essential one-way
qualification in the arithmetic application. Vanishing of $R_n(c,d)$
need not give a \emph{nonzero} common root. For example, at $d=0,c=0$,
the pair $(A_2,B_2)=(-x^2,0)$ has only the common root zero. Every
corollary above uses only the necessary implication
\[
 \text{nonzero algebraic exceptional point}
 \ \Longrightarrow\ R_n(c,d)=0.
\]

At $d=0$, one also has $A_n(0,c,0)=0$ and
\[
 B_n(0,c,0)=(-1)^{n-1}c(c+1)\cdots(c+n-2).
\]
Thus $c=0,-1,\ldots,-(n-2)$ are genuine roots of the raw resultant,
often recording an origin factor. Theorem \ref{thm:main} does not
determine all remaining integer roots, or the greatest common divisor
in $x$ at those roots.

Specifically, the argument establishes the degree conjecture in
\cite[Section 2.3(a)]{LM95}. The rational-root classifications proposed
in their parts (b),(c) remain outside the proof. The condition
$2\nu\in\Z$ includes integer orders as well as strictly half-integral
orders, and leaves infinitely many possibilities to analyze.

The closing question of \cite{LM95} for the unnormalized derivatives
$J_\nu^{(n)}$ is a different problem. Their Theorem 6 covers
$0\le n\le18$, with the nonzero algebraic exceptions
$J_{\pm1}^{(3)}(\pm\sqrt3)=0$ and
$J_0^{(4)}(\pm\sqrt3)=0$. In the present general parametrization,
the unnormalized function has $\mu=0$ and $d=-\nu$; the denominator
obstruction then gives no new restriction. No assertion about the
absence of further unnormalized exceptions is made here.

The explicit match with the 1995 conjecture is a statement about what
the proof establishes. A claim of priority over all intervening
literature would require a separate, exhaustive bibliographic review.

\begin{thebibliography}{9}
\bibitem{LM95}
L. Lorch and M. E. Muldoon,
\emph{Transcendentality of zeros of higher derivatives of functions
involving Bessel functions},
International Journal of Mathematics and Mathematical Sciences
\textbf{18} (1995), no. 3, 551--560.
\href{https://doi.org/10.1155/S0161171295000706}{doi:10.1155/S0161171295000706}.

\bibitem{Shid89}
A. B. Shidlovskii,
\emph{Transcendental Numbers}, translated by N. Koblitz,
de Gruyter Studies in Mathematics, vol. 12,
Walter de Gruyter, Berlin, 1989.
\end{thebibliography}
\end{document}
